\section{总结与展望}

\subsection{研究结论}

本文围绕“数据感知的多智能体专利分析方案自动生成”开展研究，完成了问题定义、系统设计、工程实现与实验验证四项工作。

在问题层面，本文指出现有大语言模型自动化分析系统的主要短板不在执行阶段，而在规划阶段：系统能够完成代码生成与执行，但未必能够在给定具体数据集时设计出与数据事实相匹配的高质量分析方案。该问题在专利计量场景中尤为突出，因为任务不仅要求结果产出，更要求变量解释、方法匹配与结论可复核。

在方法层面，本文提出“数据感知前置 + 双知识图谱协同 + 两轮迭代优化”的整体框架：先由 DataPreview 和 GraphPreview 提取可操作证据，再利用因果图谱与方法图谱组织假设空间和方法空间，最后通过执行反馈驱动第二轮验证性深化，从而把一次性分析升级为可追溯的探索—验证闭环。

在工程层面，本文实现了基于工作流编排框架的四智能体系统，并完成从输入研究问题到输出分析报告的端到端闭环，形成可复跑、可追踪、可审计的实验资产。

在实证层面，本文采用“主实验渐进叠加 + 组件消融 + 重复实验”的组织方式，对 6 种模式、3 个问题各重复运行 5 次，共完成 90 组实验。重复实验结果表明，完整方案 C\_iterative 在六模式总览中保持最高平均总分（23.87/25），且在相关性、针对性、方法合理性、创新性和深度五个维度上均表现稳定。结合实验结果，可以形成以下四项具体结论：

第一，数据感知是提升方案质量的稳定关键。与 A 和 D 相比，数据感知模式 B 在重复实验中稳定取得更高的语义评分和更强的数据证据引用，说明在规划前显式引入数据分布、图结构与样本语义证据，能够有效推动方案从“问题语义驱动”转向“数据证据驱动”。

第二，两轮迭代是拉开系统上限的核心机制。C\_iterative 在主实验四模式中显著高于其他方案，不仅总分最高，而且在方法多样性、参数丰富度、任务数和五维语义评分上均呈现更均衡的领先，说明第二轮并非简单重复第一轮，而是围绕证据缺口进行补证据、换方法和做稳健性分析。

第三，知识图谱具有正向但任务依赖的贡献。组件消融结果显示，移除知识图谱后总体平均表现由 C 的 23.87 降至 F 的 22.20，说明知识图谱在平均意义上仍提供正向增益；但分问题结果同时表明，图谱收益并不均匀，在结构锚点清晰、分析链条较长的问题上更明显，在开放度更高的问题上则需要与数据证据更紧密协同，避免静态路径过度约束。

第四，系统已表现出较好的工程可用性。6 种模式在 90 组实验中的代码执行成功率均为 100\%，表明当前系统在所测试的三个研究问题上能够稳定完成规划、执行与结果落盘，实验差异主要来自规划质量而非执行失败。

在研究贡献方面，理论上，本文将自动化分析研究的关注点从“执行正确性”前移到“规划合理性”，并进一步把知识图谱的作用讨论从“是否有效”细化为“在什么条件下更有效”。方法上，本文给出可复用的系统范式：数据感知组件负责证据抽取，双图谱组件负责结构约束，迭代机制负责深度提升，三者组合可迁移到其他结构化分析场景。实践上，本文形成了可落地工程资产，包括主流程代码、90 组实验输出和评估脚本，可直接支持后续复验与扩展开发。

\subsection{研究展望}

尽管本文已完成 90 组重复实验，当前研究仍存在以下局限：

第一，统计推断仍不充分。本文已经通过每种模式在每个问题上的 5 次独立重复降低了单次运行偏差，但尚未进一步引入 Friedman 检验、Wilcoxon 符号秩检验或 bootstrap 置信区间对关键差异进行正式显著性分析，因此当前结论更适合理解为“具有重复证据支持的趋势判断”。

第二，评估主体仍然单一。LLM-as-Judge 虽采用匿名并列评分与低温度参数缓解偏差，但仍可能存在对表述风格、篇幅长度或特定方法的系统性偏好，当前尚未加入人工专家盲审进行交叉校验。

第三，问题规模与外部有效性仍有限。本文仅覆盖三个研究问题和单一领域数据集，虽已能够观察到数据感知、知识图谱和迭代机制的差异化作用，但能否外推到更多专利分析任务、更多行业领域和更多数据来源，仍需进一步验证。

第四，知识约束的自适应能力仍有改进空间。当前图谱仍以静态文献归纳结果为主，尚未根据数据证据强度、问题结构差异和迭代反馈动态调节约束强度，因此在开放问题上的收益仍不稳定。

第五，跨模型与跨环境可迁移性仍待验证。本文控制了统一模型版本、提示模板和执行环境，但在不同模型、不同参数和不同计算环境下是否仍保持类似的相对排序，仍需后续补充。

围绕上述局限，后续研究可从以下五个方向继续推进：

\begin{enumerate}[(1)]
\item 在已有 $n=5$ 重复实验的基础上补充正式统计检验。针对关键对比（如 B vs D、C vs B、C vs F）计算 bootstrap 置信区间，并引入 Friedman 检验或 Wilcoxon 符号秩检验评估差异显著性。
\item 构建“自动评估 + 专家复核”的混合评审框架。将 LLM-as-Judge 与领域专家盲审结合，并引入评分者一致性指标（如 Krippendorff's $\alpha$）校准大语言模型评分的可靠性。
\item 引入自适应约束强度机制。根据当前数据洞察的信号强度、问题结构类型和第一轮执行反馈动态调节图谱约束权重，使知识图谱从静态模板升级为可被数据重塑的引导器。
\item 优化数据感知的广度—深度平衡。针对不同问题类型设计动态任务数校准策略，在保持高针对性的同时增加最低覆盖率约束，避免系统在开放问题中过早收缩探索空间。
\item 推进跨领域迁移验证。将本文框架扩展到生物医药、新能源和半导体等技术领域，检验数据感知机制、双图谱框架和迭代闭环在不同领域中的可迁移性与边界条件。
\end{enumerate}

总体而言，本文为数据感知多智能体专利分析框架提供了较为扎实的重复实验支持，并揭示了数据感知、知识图谱和迭代机制在专利分析规划任务中的差异化作用。这些发现不仅为当前系统的继续优化提供了方向，也为后续面向更大规模、更复杂场景的自动化科学分析系统奠定了方法与工程基础。
